% !PRINTED 132
\weberpage{132}
%<<CONT
a function which, by § 36, remains completely unchanged under linear transformation. If one makes further use of the relation (2):
\[
\eta_1'\omega_1 - \eta_1\omega_1' = -\pi i \beta,
\]
then it follows that
\weq{8}{
\frac{\vartheta_{00}(v,\omega)\vartheta_{01}(v,\omega)\vartheta_{10}(v,\omega)}{\vartheta_{00}\vartheta_{01}\vartheta_{10}}
= e^{-3\pi i \beta v v'}\, \frac{\vartheta_{00}(v',\omega')\vartheta_{01}(v',\omega')\vartheta_{10}(v',\omega')}{\vartheta_{00}(0,\omega')\vartheta_{01}(0,\omega')\vartheta_{10}(0,\omega')},
}
in which one may still, by (19) of the preceding section,
\[
\vartheta_{00}(0,\omega')\vartheta_{01}(0,\omega')\vartheta_{10}(0,\omega') = \varepsilon^3\sqrt{\alpha+\beta\omega}^{\,3}\,\vartheta_{00}\vartheta_{01}\vartheta_{10}
\]
set.

But we draw another conclusion from (8): namely, if one puts
\[
v' = \frac{h}{n}, \quad v = \frac{h(\alpha+\beta\omega)}{n},
\]
and takes the product for $h=1,2,\cdots,\dfrac{n-1}{2}$, then there results, by means of the known relation (Bd.\ I, § 11)
\[
\sum_{1,\ \frac{n-1}{2}}^{h} h^2 = n\,\frac{n^2-1}{24},
\]
if one applies formula (23), § 32 to the right-hand side of (8):
\weq{9}{
\begin{aligned}
&2^{\frac{n-1}{2}} \prod_{1,\ \frac{n-1}{2}}^{h} \vartheta_{00}\!\left(\frac{h(\alpha+\beta\omega)}{n}\right)\vartheta_{01}\!\left(\frac{h(\alpha+\beta\omega)}{n}\right)\vartheta_{10}\!\left(\frac{h(\alpha+\beta\omega)}{n}\right) \\
&= e^{-\frac{\pi i}{8}\frac{n^2-1}{n}\beta(\alpha+\beta\omega)}\ \vartheta_{00}^{\frac{n-1}{2}}\,\vartheta_{01}^{\frac{n-1}{2}}\,\vartheta_{01}^{\frac{n-1}{2}}, % [sic] German source prints \vartheta_{01} twice; the third factor is mathematically expected to be \vartheta_{10}.
\end{aligned}
}
in which now $\alpha, \beta$ may be any pair of relatively prime numbers.

\begin{center}
\textbf{§ 40. \gesperrt{Linear Transformation of the Functions}\\
$f(\omega), f_1(\omega), f_2(\omega)$.}
\end{center}

By § 34, (9), the formulae for the linear transformation of the $f$-functions are a simple consequence of the transformation of the $\eta$-function.

If we set in the linear transformation
\[
\begin{pmatrix}\alpha, & \beta\\ \gamma, & \delta\end{pmatrix}
\]
%CONT>>
